\section{USAMO 1986}
        \begin{enumerate}
        	\item (\textit{USAMO 1986}) $(\text{a})$ Does there exist 14 consecutive positive integers each of which is divisible by one or more primes $p$ from the interval $2\le p \le 11$?


 $(\text{b})$ Does there exist 21 consecutive positive integers each of which is divisible by one or more primes $p$ from the interval $2\le p \le 13$?


 



\textbf{Solution}

(a)     To solve part (a), we first note that for any 14 consecutive positive integers, exactly 7 are even (divisible by 2) and therefore satisfy the criteria. We can remove these from the problem, and simplify it to the following question, which is equivalent to part (a):"Do there exist 7 consecutive positive odd integers each of which is divisible by one or more primes $p$ from the interval $3\le p \le 11$?"Among any 7 consecutive positive odd integers, the following holds:
 \[\text{Either 2 or 3 are divisible by 3}\]
 \[\text{Either 1 or 2 are divisible by 5}\]
 \[\text{Exactly 1 is divisible by 7}\]
 \[\text{Either 0 or 1 are divisible by 11}\]For every one of these seven integers to be divisible by one of 3, 5, 7, or 11, there must be 3 multiples of 3, 2 multiples of 5, 1 multiple of 7, and 1 multiple of 11. Additionally, none of these integers may be a multiple of any two of the four aforementioned primes. Otherwise, the Pigeonhole Principle dictates that at least one integer is not divisible by any of the four, thus failing to meet the criteria. But this cannot be. Calling the consecutive odd integers $a_1, a_2, \dots, a_7$, we note that $a_1$, $a_4$, and $a_7$ must be multiples of 3, and therefore cannot be multiples of 5, 7, or 11. But for there to be two multiples of 5, they must be one of the two pairs $(a_1,a_6)$ and $(a_2,a_7)$. But each of these pairs contains a multiple of 3, and so at least one of the 7 odd integers is divisible by none of the primes 3, 5, 7, or 11. Therefore the answer to part (a) is no.

 (b)    To solve part (b), we use a strategy similar to the one used in part (a), reducing part (b) to this equivalent question:"Do there exist 10 consecutive positive odd integers each of which is divisible by one or more primes $p$ from the interval $3\le p \le 13$?"(We will ignore the case where the first of the 21 integers is odd, resulting in 11 odd integers instead of 10, as it is only true if the weaker, 10-integer argument is also true.)Among any 10 consecutive positive odd integers, the following holds:
 \[\text{Either 3 or 4 are divisible by 3}\]
 \[\text{Exactly 2 are divisible by 5}\]
 \[\text{Either 1 or 2 are divisible by 7}\]
 \[\text{Either 0 or 1 are divisible by 11}\]
 \[\text{Either 0 or 1 are divisible by 13}\]For every one of these ten integers to be divisible by one of 3, 5, 7, 11, or 13, there must be 4 multiples of 3, 2 multiples of 5, 2 multiples of 7, 1 multiple of 11, and 1 multiple of 13. As before, none of these integers may be a multiple of any two of the five primes. Calling the consecutive odd integers $a_1, a_2, \dots, a_{10}$, we note that $a_1$, $a_4$, $a_7$, and $a_{10}$ must be multiples of 3, and therefore cannot be multiples of 5, 7, 11, or 13. Unlike part (a), however, this stipulation is not a dead end. Let the multiples of 5 be $a_3$ and $a_8$ and let the multiples of 7 be $a_2$ and $a_9$. The multiples of 11 and 13 are $a_5$ and $a_6$, in some order (it doesn't really matter which). An example of a sequence of 21 consecutive positive integers satisfying part (b) is the integers from 9440 to 9460 (inclusive), which can be obtained by solving modular equations that result from these statements. So the answer to part (b) is yes.

 


	\item (\textit{USAMO 1986}) During a certain lecture, each of five mathematicians fell asleep exactly twice. For each pair of mathematicians, there was some moment when both were asleep simultaneously. Prove that, at some moment, three of them were sleeping simultaneously.


 



\textbf{Solution}

We shall assume to the contrary that there was never a time when three mathematicians were sleeping simultaneously, and derive a contradiction.

 As a subtle but logically necessary note, we will assume without loss of generality that no two events (an event is one mathematician either falling asleep or waking up) happen at the same time.

 
            \begin{enumerate}[label=(\roman*)]
                \item Rationale: Consider a group of events happening simultaneously--i.e., there is a moment before any of the events have happened, and there is a moment after all the events have happened, but there are no moments in which at least one but not all of the events have happened.  If there is a valid timeline--i.e. one in which every pair of mathematicians was simultaneously asleep at some moment, but in which no three mathematicians are simultaneously asleep at any moment--containing a group of simultaneous events, then there is also a valid timeline in which that group of events is broken up into a sequence; you can construct the second timeline by putting the "mathematician wakes up" events first and the "mathematician falls asleep" events last.  (In the first timeline, there were n mathematicians asleep in the moments immediately before the group of events, and m mathematicians asleep immediately afterwards; in the second timeline, there will be additional moments when n-1, n-2, ..., k, k+1, ..., m-1 mathematicians were asleep.  If the first timeline is valid, then n and m are both less than 3, so the additional intermediate moments will not violate the "no 3 asleep at once" rule; and since we did not remove any moments, if the "all sleeping-pairs" rule was satisfied in the first timeline, then it is satisfied in the second timeline.)  So the existence of any valid timeline implies the existence of a sequential valid timeline.
            \end{enumerate}
            We have (5 choose 2) = 10 sleeping-pairs of mathematicians to account for.  Also, since 5 mathematicians each fell asleep 2 times, we have a total of 10 occasions on which a mathematician fell asleep.  Now, if two mathematicians are ever asleep simultaneously, then one of them fell asleep while the other was already asleep.  However, because of the "no three asleep at once" rule, if a mathematician falls asleep, then at most one other mathematician could have been asleep already.  Therefore, each occasion when a mathematician falls asleep can account for at \textit{most} one sleeping-pair.  It would appear that we have just enough to make it.  However, when the first mathematician falls asleep, no other mathematicians are asleep...

 Hey, wait a minute.  Couldn't one mathematician, or even two, have been already asleep when the lecture started?  In fact, the wording of the problem does not forbid this, and this permits us to construct an easy counterexample:

 Call the mathematicians A,B,C,D,E.  A starts out asleep.  B sleeps; A wakes; C sleeps; B wakes; D sleeps; C wakes; E sleeps; D wakes; A sleeps; E wakes; then, C sleeps; A wakes; E sleeps; C wakes; B sleeps; E wakes; D sleeps; B wakes; A sleeps; D wakes.  This creates, in order, the sleeping-pairs AB, BC, CD, DE, EA, then AC, CE, EB, BD, DA; each mathematician falls asleep and wakes up exactly twice, and at no time are three mathematicians asleep.

 Well, this must be considered a hole in the problem as written.  However, if we add the (nontrivial) assumption that the mathematicians were all awake when the lecture began, then it follows that the first occasion when a mathematician falls asleep can only account for zero sleeping-pairs, and the remaining 9 occasions can only account for 9 of 10 sleeping-pairs, and so there must be some pair of mathematicians that are not simultaneously asleep during the lecture.  So our assumption to the contrary implies an impossible sequence of events.  Therefore, that assumption must be wrong, and there must be a moment when three mathematicians are asleep.

 


	\item (\textit{USAMO 1986}) What is the smallest integer $n$, greater than one, for which the root-mean-square of the first $n$ positive integers is an integer?


 $\mathbf{Note.}$ The root-mean-square of $n$ numbers $a_1, a_2, \cdots, a_n$ is defined to be

 \[\left[\frac{a_1^2 + a_2^2 + \cdots + a_n^2}n\right]^{1/2}\]


 



\textbf{Solution 1}

We use the fact that $\sum_{i=1}^n i^2=\frac{n(n+1)(2n+1)}{6}$.Thus the root mean square is an integer if and only if $\frac{n(n+1)(2n+1)}{6n}=\frac{(n+1)(2n+1)}{6}$ is an integer. Since $\gcd(n+1,2n+1)=1$ and $2n+1$ is not even, $n+1$ is even, so $n$ is odd. Thus, $2n+1\equiv3\pmod4$. Since perfect squares are only $0$ or $1$ modulo $4$, $2n+1$ cannot be a perfect square. Therefore, there are nonnegative integers $a$ and $b$ with $n+1=2a^2$ and $2n+1=3b^2$, which means that $n\equiv1\pmod3$. Since $b$ is odd, $b^2\equiv1\pmod8\Rightarrow n\equiv1\pmod4$. Thus $a$ must be an odd integer not divisible by $3$. We test $a=5$, $7$, $11$, all of which do not work. Testing $a=13$ gives $(n,b)=(337,15)$, therefore our answer is $\boxed{337}$.

 



\textbf{Solution 2}

Let's first obtain an algebraic expression for the root mean square of the first $n$ integers, which we denote $I_n$. By repeatedly using the identity$(x+1)^3 = x^3 + 3x^2 + 3x + 1$, we can write
 \[1^3 + 3\cdot 1^2 + 3 \cdot 1 + 1 = 2^3,\]
 \[1^3 + 3 \cdot(1^2 + 2^2) + 3 \cdot (1 + 2) + 1 + 1  = 3^3,\] and
 \[1^3 + 3\cdot(1^2 + 2^2 + 3^2) + 3 \cdot (1 + 2 + 3) + 1 + 1 + 1 = 4^3.\] We can continue this pattern indefinitely, and thus for anypositive integer $n$,
 \[1 + 3\sum_{j=1}^n j^2 + 3 \sum_{j=1}^n j^1 + \sum_{j=1}^n j^0 = (n+1)^3.\]Since $\sum_{j=1}^n j  = n(n+1)/2$, we obtain
 \[\sum_{j=1}^n j^2 = \frac{2n^3 + 3n^2 + n}{6}.\]Therefore,
 \[I_n = \left(\frac{1}{n} \sum_{j=1}^n j^2\right)^{1/2} = \left(\frac{2n^2 + 3n + 1}{6}\right)^{1/2}.\]Requiring that $I_n$ be an integer, we find that
 \[(2n+1 ) (n+1) = 6k^2,\] where $k$ is an integer. Using the Euclidean algorithm, we see that$\gcd(2n+1, n+1) = \gcd(n+1,n) = 1$, and so $2n+1$ and $n+1$ share nofactors greater than 1. The equation above thus implies that $2n+1$and $n+1$ is each proportional to a perfect square. Since $2n+1$ isodd, there are only two possible cases:

 Case 1: $2n+1 = 3 a^2$ and $n+1 = 2b^2$, where $a$ and $b$ are integers.

 Case 2: $2n+1 = a^2$ and $n+1 = 6b^2$.

 In Case 1,  $2n+1 = 4b^2 -1 = 3a^2$. This means that $(4b^2 -1)/3 = a^2$ for some integers $a$ and $b$. We proceed by checking whether$(4b^2-1)/3$ is a perfect square for $b=2, 3, 4, \dots$. (The solution$b=1$ leads to $n=1$, and we are asked to find a value of $n$ greater than 1.) The smallest positive integer $b$ greater than 1 forwhich $(4b^2-1)/3$ is a perfect square is $b=13$, which results in $n=337$.

 In Case 2, $2n+1 = 12b^2 - 1 = a^2$. Note that $a^2 = 2n+1$ is an odd square, and hence is congruent to $1 \pmod 4$. But $12b^2 -1 \equiv 3 \pmod 4$ for any $b$, so Case 2 has no solutions.

 Alternatively, one can proceed by checking whether $12b^2 -1$ is a perfect square for $b=1, 2 ,3 ,\dots$. We find that $12b^2 -1$ is not a perfect square for $b = 1,2, 3, ..., 7, 8$, and $n= 383$ when $b=8$. Thus the smallest positive integers $a$ and $b$ for which $12b^2- 1 = a^2$ result in a value of $n$ exceeding the value found in Case 1, which was 337.

 In summary, the smallest value of $n$ greater than 1 for which $I_n$ is an integer is $\boxed{337}$.

 


	\item (\textit{USAMO 1986}) Two distinct circles $K_1$ and $K_2$ are drawn in the plane. They intersect at points $A$ and $B$, where $AB$ is the diameter of $K_1$. A point $P$ on $K_2$ and inside $K_1$ is also given.


 Using only a "T-square" (i.e. an instrument which can produce a straight line joining two points and the perpendicular to a line through a point on or off the line), find a construction for two points $C$ and $D$ on $K_1$ such that $CD$ is perpendicular to $AB$ and $\angle CPD$ is a right angle.


 



\textbf{Solution}

<i>This problem needs a solution. If you have a solution for it, please help us out by \href{https://artofproblemsolving.com/wiki/index.php?title=1986\_USAMO\_Problems/Problem\_4&action=edit}{adding it}.</i> 

 


	\item (\textit{USAMO 1986}) By a partition $\pi$ of an integer $n\ge 1$, we mean here a representation of $n$ as a sum of one or more positive integers where the summands must be put in nondecreasing order. (E.g., if $n=4$, then the partitions $\pi$ are $1+1+1+1$, $1+1+2$, $1+3, 2+2$, and $4$).


 For any partition $\pi$, define $A(\pi)$ to be the number of $1$'s which appear in $\pi$, and define $B(\pi)$ to be the number of distinct integers which appear in $\pi$. (E.g., if $n=13$ and $\pi$ is the partition $1+1+2+2+2+5$, then $A(\pi)=2$ and $B(\pi) = 3$).


 Prove that, for any fixed $n$, the sum of $A(\pi)$ over all partitions of $\pi$ of $n$ is equal to the sum of $B(\pi)$ over all partitions of $\pi$ of $n$.


 



\textbf{Solution}

Let $S(n) = \sum\limits_{\pi} A(\pi)$ and let $T(n) = \sum\limits_{\pi} B(\pi).$ We will use generating functions to approach this problem -- specifically, we will show that the generating functions of $S(n)$ and $T(n)$ are equal.

 Let us start by finding the generating function of $S(n).$ This function counts the total number of 1's in all the partitions of $n.$ Another way to count this is by counting the number of partitions of $n$ that contain $x$ 1's and multiplying this by $x,$ then summing for $1\leq x \leq n.$ However, the number of partitions of $n$ that contain $x$ 1's is the same as the number of partitions of $n-x$ that contain no 1's, so

 
 \[S(n) = \sum_{k=1}^n k\cdot (\text{\# of partitions of }n-k\text{ with no 1's}).\]

 The number of partitions of $m$ with no 1's is the coefficient of $x^m$ in 

 
 \begin{align*} F(x) &= (1+x^2+x^4+\ldots)(1+x^3+x^6+\ldots)(1+x^4+x^8+\ldots)\ldots \\ &= \prod\limits_{i=2}^{\infty}\frac{1}{1-x^i} \end{align*}


 Note that there is no $(1+x+x^2+x^3+\ldots)$ term in $F(x)$ because we cannot have any 1's in the partition.

 Let $c_m$ be the coefficient of $x^m$ in the expansion of $F(x),$ so we can rewrite it as $F(x)=c_0+c_1x+c_2x^2+c_3x^3+\ldots.$ We wish to compute $S(n)=1\cdot c_{n-1}+2\cdot c_{n-2}+\ldots+n\cdot c_0.$ 

 Consider the power series $G(x)=(x+2x^2+3x^3+4x^4+\ldots)F(x).$

 
 \begin{align*} G(x) &= (x+2x^2+3x^3+4x^4+\ldots)(c_0+c_1x+c_2x^2+c_3x^3+\ldots)\\ &= x(1+2x+3x^2+4x^3+\ldots) \prod\limits_{i=2}^{\infty}\frac{1}{1-x^i}\\ &= x\cdot\frac{1}{(1-x)^2} \prod\limits_{i=2}^{\infty}\frac{1}{1-x^i}\\ &= \frac{x}{1-x} \prod\limits_{i=1}^{\infty}\frac{1}{1-x^i} \end{align*}


 If we expand the first line, we see that the coefficient of $x^n$ in $G(x)$ for any $n$ is $1\cdot c_{n-1}+2\cdot c_{n-2}+\ldots+n\cdot c_0,$ which is exactly $S(n)$! So by definition, $G(x)=\frac{x}{1-x} \prod\limits_{i=1}^{\infty}\frac{1}{1-x^i}$ is the generating function of $S(n).$

 Now let's find the generating function of $T(n).$ Notice that counting the number of distinct elements in each partition and summing over all partitions is equivalent to counting how many partitions of $n$ contain i and then summing for all $1\leq i \leq n.$ 

 So, 

 
 \[T(n) = \sum_{k=1}^n k\cdot (\text{\# of partitions of }n\text{ that contain a }k).\]

 However, the number of partitions of n that contain a $i$ is the same as the total number of partitions of $n-i,$ so 
 \[T(n) = \sum_{k=1}^n k\cdot (\text{\# of partitions of }n-k).\]

 The generating function for the number of partitions is 
 \begin{align*} P(x) &= (1+x+x^2+\ldots)(1+x^2+x^4+\ldots)(1+x^3+x^6+\ldots)\ldots \\ &= \prod\limits_{i=1}^{\infty} \frac{1}{1-x^i}. \end{align*}


 Let's write the expansion of $P(x)$ as $P(x)=d_0+d_1x+d_2x^2+d_3x^3+\ldots,$ so we wish to find $T(n)=d_0+d_1+d_2+\ldots+d_{n-1}.$

 Consider the power series $H(x)=(x+x^2+x^3+x^4+\ldots)P(x).$

 
 \begin{align*} H(x) &= (x+x^2+x^3+x^4+\ldots)(d_0+d_1x+d_2x^2+d_3x^3+\ldots) \\ &= x(1+x+x^2+x^3+\ldots) \prod\limits_{i=1}^{\infty} \frac{1}{1-x^i}\\ &= x\cdot \frac{1}{1-x}\prod\limits_{i=1}^{\infty} \frac{1}{1-x^i}\\ &= \frac{x}{1-x} \prod\limits_{i=1}^{\infty} \frac{1}{1-x^i} \end{align*}


 If we expand the first line, we see that the coefficient of $x^n$ in $H(x)$ for any $n$ is $d_{n-1}+d_{n-2}+\ldots+d_0,$ which is precisely $T(n).$ This means $H(x)$ is the generating function of $T(n).$

 Thus, the generating functions of $S(n)$ and $T(n)$ are the same, so $S(n)=T(n)$ for all $n$ and we are done.

 



\textbf{Solution 2}

First of all, the generating function for all partitions of $n$ is 
 \[\frac{1}{(1-x)(1-x^2)(1-x^3)\cdots}\]Let us define 
 \[s(n) = \{\text{\# of }\Pi\text{ for }n\text{ with one 1 in the partition}\} + 2\cdot\{\text{\# of }\Pi\text{ for }n\text{ with two 1s}\} + \cdots\]Now, we make a great discovery. 
 \[\{\text{\# of }\Pi\text{ with for }n\text{ with one 1}\} = \{\text{\# of }\Pi\text{ for }n-1\text{ with no 1's}\}\]This is because when $n$ has one $1$ in a partition, then if we remove that $1$, then we get the number $n-1$, and our partition no longer has any $1$s.

 So, we can rewrite $s(n)$ as 
 \[s(n) = \{\text{\# of }\Pi\text{ for }n-1\text{ with no 1 in the partition}\} + 2\cdot\{\text{\# of }\Pi\text{ for }n-2\text{ with no 1s}\} + \cdots\]Let us define 
 \[a_n = \{\text{\# of }\Pi\text{ for }n\text{ with no 1s}\}\]Then, 
 \[s(n) = a_{n-1} + 2a_{n-2} + \cdots\]

 The generating function for $a_n$ (call it $\alpha(n)$) is just 
 \[\frac{1}{(1-x^2)(1-x^3)\cdots}\]We removed the $(1-x)$ because $\Pi \in a_n$ cannot have any $1$s.Now, let us try to find the generating function for $s(n)$.

 We also know that 
 \[\alpha(n) = a_1x + a_2x^2 + \cdots\]By the definition of $s(n)$, we know that its generating function (call it $\sigma(n)$) looks something like this:
 \[\sigma(n) = a_1x^2+(2a_1+a_2)x^3+(3a_1+2a_2+a_3)x^4 + \cdots\]Let's try to write this in terms of $\alpha(n)$. We know that $x\alpha(n)$ is 
 \[x\alpha(n) = a_1x^2+a_2x^3 + \cdots\]Now, when we do $\sigma(n) -x \alpha(n)$, we will get that $a_1x^2, a_2x^3, a_3x^4, \cdots$ cancel out.

 However, we still need more to cancel out.

 Let's focus on that $2a_1x^3$ that didn't cancel out and is in $\sigma(n)$.To make that cancel out, we'd have to do $\sigma(n) - x\alpha(n) - 2x^2\alpha(n)$, where 
 \[2x^2\alpha(n) = 2a_1x^3+2a_2x^4 + \cdots\]See the $2a_1x^3$ in both $\sigma(n)$ and $2x^2\alpha(n)$?

 Great, now we focus on $3a_1x^4$. Just like last time, we'd have to do $\sigma(n) - x\alpha(n) - 2x^2\alpha(n) - 3x^3\alpha(n)$, where 
 \[3x^3\alpha(n) = 3a_1x^4+3a_2x^5 + \cdots\]

 In order to make all of that stuff in $\sigma(n)$ cancel out, following that pattern, we would do 
 \[\sigma(n) - \alpha(n)\cdot(x+2x^2+3x^3+4x^4 + \cdots) = 0\]or 
 \[\sigma(n) = \alpha(n) \cdot (x+2x^2+3x^3+4x^4 + \cdots)\]So what's a better way to write $x+2x^2+3x^3+4x^4+\cdots$? 

 Well, we already know that 
 \[\frac{x}{1-x} = x+x^2+x^3 + \cdots\]And also 
 \[\frac{x^2}{1-x} = x^2+x^3+x^4+\cdots\]So adding these up:
 \[\begin{array}{rrrrrrr} & x & +x^2 & +x^3 & +x^4 & +x^5 & +\cdots \\ +&   & x^2 & +x^3 & +x^4 & +x^5 & +\cdots \\ \hline & x & +2x^2 & +2x^3 & +2x^4 & +2x^5 & +\cdots \end{array}\]We already have the first few terms matching! Now, if we add on $\dfrac{x^3}{1-x} = x^3+x^4+x^5+\cdots$, we get 
 \[\begin{array}{rrrrrrrr} & x & +x^2 & +x^3 & +x^4 & +x^5 & +x^6 & +\cdots \\ +&   & x^2 & +x^3 & +x^4 & +x^5 & +x^6 & +\cdots \\ +&   &     & x^3 & +x^4 & +x^5 & +x^6 & +\cdots \\ \hline & x & +2x^2 & +3x^3 & +3x^4 & +3x^5 & +3x^6 & +\cdots \end{array}\]We see that in order to get our wanted polynomial, following the pattern, we'll have to do 
 \[x+2x^2+3x^3+4x^4+\cdots = \frac{x}{1-x} + \frac{x^2}{1-x}+\frac{x^3}{1-x}+\cdots\]However, this can be written as the infinite geometric series with $a=\frac{x}{1-x}$ and $r=x$.So, 
 \[x+2x^2+3x^3+\cdots = \frac{\frac{x}{1-x}}{1-x}\]Thus, 
 \[\sigma(n) = \frac{1}{(1-x^2)(1-x^3)\cdots} [\text{because of }\alpha(n)] \cdot \frac{x}{(1-x)^2} [\text{because of the thing we found above}]\]

 Now, we go to $t(n)$. $t(n)$ is $\sum_{\Pi \text{ is a partition}} g(\Pi)$, where $g(\Pi)$ is the number of distinct numbers in a partition.

 Let $r(1,n)$ be the number of $\Pi$ that have a $1$ in them$r(2, n)$ be the number of $\Pi$ that have a $2$ in them...

 Then, 
 \[t(n) = r(1, n) + r(2,n) + \cdots + r(n,n)\]by the definition of $t(n)$.

 Now, we notice that the generating function for $r(1,n)$ is $x \cdot \text{normal partition gen function}$. because $r(1,n)$ is the number of partitions of $n-1$ ($p_{n-1}$) (it is a normal partion of $n-1$ with a $1$ added on).$r(2,n)$ is the number of partitions of $n-2$ ($p_{n-2}$)

 So now, we see that 
 \[t(n) = p_{n-1} + p_{n-2} + \cdots\]We want to write down the generating function for $t(n)$, call it $\tau(x)$.When $n=2$:$t(2)=p_1$$t(3)=p_1+p_2$

 Similarly, by choosing some values of $n$, we get 
 \[\tau(x) = p_1x^2 + (p_1+p_2)x^3+\cdots\]If $\rho(x)$ is the partition generating function, then 
 \[\rho(x) = p_1x+p_2x^2 + \cdots\]And, in that case, we can substitute in $\rho(x)$, into $\tau(x)$, to get 
 \[\tau(x) = \rho(x)(x+x^2+\cdots)\]The second term in this expression ($x+x^2+\cdots$) can be written as the infinite sum $\frac{x}{1-x}$, so we can simplify: 
 \[\tau(x)=\frac{1}{(1-x)(1-x^2)\cdots} \cdot \frac{x}{1-x}\]

 From above, we found that 
 \[\sigma(x) = \frac{1}{(1-x^2)(1-x^3)\cdots} \cdot \frac{x}{(1-x)^2}\]If we can prove that $\tau(x) = \sigma(x)$, then we are done!They are equal, so we have proved the given statement.

 Q.E.D

 


\end{enumerate}